%% ─── Related Work ────────────────────────────────────────────────────────────
\section{Related Work}
\label{sec:related}

\subsection{LLM-Based Code Agents and Tools}

SWE-Agent~\cite{yang2024sweagent} and Agentless~\cite{xia2024agentless} are
representative systems in which agents navigate repositories through file-read
and string-edit operations.  AutoCodeRover~\cite{zhang2024autocoderover}
augments this with symbol indices and file maps to guide localization.
CODESTRUCT~\cite{kim2025codestruct} replaces text-based edits with AST-grounded
primitives (\texttt{readCode}/\texttt{editCode}), reducing input tokens by
12--38\% and improving SWE-Bench Verified Pass@1 by up to 5.0pp.  KB is
complementary to these systems: whereas CODESTRUCT improves the \emph{action
interface} once a target is located, KB improves \emph{context retrieval} by
replacing filesystem exploration with structured fact lookup.  The two
approaches can be composed.

\subsection{Retrieval-Augmented Generation}

RAG~\cite{lewis2020rag} grounds LLM answers in retrieved passages, but standard
RAG pipelines are designed for static document collections.  They do not model
the inter-entity relationships that characterize codebases (function calls,
class hierarchies, module imports), and they treat retrieval as a one-shot
lookup rather than an adaptive loop.  KB's multi-pond BFS loop merges five
evidence streams per pass and exits on a scored sufficiency condition, making
retrieval quality a function of graph density rather than fixed index size.

\subsection{LLM Evaluation Frameworks}

Evaluating LLM outputs with another LLM (LLM-as-judge) is established in
MT-Bench~\cite{zheng2023judging} and AlpacaFarm~\cite{dubois2024alpacafarm}.
Known failure modes include position bias, verbosity bias, and self-enhancement
bias.  PandaLM~\cite{wang2023pandalm} trains a dedicated judge model to
mitigate these.  MOEL's jury module addresses the same biases at inference time
through position debiasing (swapped-order consistency), verbosity normalization
($\log_{1+}$ length scaling), self-enhancement down-weighting ($\times 0.5$ for
same-family judges), and family diversity enforcement---without requiring a
fine-tuned judge.  Unlike prior work, MOEL is not outcome-only: it combines
the LLM jury with an AST-structural loss and trajectory/resource losses so that
efficiency failures are numerically visible.

\subsection{Agent Efficiency Measurement}

SWE-ContextBench measures token cost and cache efficiency across context
injection strategies; CodeScaleBench evaluates tool validity across repository
scale classes.  Neither framework provides a unified scalar loss that spans
correctness, trajectory, and resource consumption simultaneously.  The closest
prior work is SWE-Atlas, which verifies agent-declared file manifests against
live \texttt{git diff} output.  MOEL's \texttt{ManifestValidator} adopts the
same live-diff approach and extends it with a \texttt{MutationValidator} that
applies tree-sitter AST body stubs to confirm causal linkage between agent
changes and test coverage.

\subsection{AST-Based Code Analysis}

Tree-sitter~\cite{falleri2014gumtree} and related tools provide fine-grained
AST differencing.  ASTNN~\cite{alon2019code2vec} uses AST path embeddings for
code representation learning.  MOEL's $L_\text{AST}$ adopts a simpler but
interpretable measure: Jaccard distance over named exported declarations,
computed deterministically without learned components, making it reproducible
across model versions.
